\documentclass[12pt,a4paper]{article}
\usepackage[margin=1in]{geometry}
\usepackage{amsmath}
\usepackage{amssymb}
\usepackage{array}
\usepackage{booktabs}
\usepackage{xcolor}
\usepackage{fancyhdr}
\usepackage{graphicx}
\usepackage{setspace}

\onehalfspacing
\pagestyle{fancy}
\fancyhf{}
\fancyhead[L]{CPE 415 -- Digital Image Processing}
\fancyhead[R]{Mid Term SP25}
\fancyfoot[C]{\thepage}
\renewcommand{\headrulewidth}{0.5pt}

\title{\textbf{Mid Term Examination -- Spring 2025}}
\author{CPE 415: Digital Image Processing}
\date{}

\begin{document}

\section*{Examination Details}
\vspace{-0.5cm}
\noindent\begin{tabular}{l@{\hspace{3cm}}l}
\textbf{Course Title:} & Digital Image Processing (CPE 415) \\
\textbf{Credit Hours:} & 4(3,1) \\
\textbf{Instructor:} & Mr. Moazzam Ali \\
\textbf{Program:} & Bachelor of Engineering, Computer Engineering \\
\textbf{Semester:} & Spring 2025 \\
\textbf{Total Marks:} & 50 \\
\textbf{Time Allowed:} & 90 Minutes \\
\textbf{Institution:} & COMSATS University Islamabad, Lahore Campus
\end{tabular}

\vspace{1.0cm}
\noindent\textbf{Examination Instructions:}
\begin{enumerate}
\item This is a closed-book examination.
\item Use of mobile phones and calculators is strictly prohibited during the examination.
\item Candidates must attempt all questions; no choice is permitted.
\item All answers must be written in the provided answer booklet.
\end{enumerate}

\vspace{1.0cm}
\section*{Question 1 \quad [CLO1-PL01-C4] \quad 10 Marks}

A $7 \times 10$ 4-bit grayscale image is provided and represented by the following matrix $T(I)$:

\begin{equation*}
T(I) = \begin{pmatrix}
0 & 0 & 2 & 6 & 8 & 4 & 0 & 0 & 0 & 0 \\
2 & 4 & 6 & 15 & 15 & 8 & 4 & 2 & 0 & 0 \\
6 & 8 & 15 & 15 & 15 & 15 & 8 & 4 & 0 & 0 \\
4 & 6 & 8 & 15 & 15 & 15 & 6 & 2 & 0 & 0 \\
2 & 4 & 6 & 15 & 15 & 8 & 4 & 0 & 0 & 0 \\
0 & 2 & 4 & 8 & 15 & 6 & 2 & 0 & 0 & 0 \\
0 & 0 & 2 & 6 & 8 & 4 & 0 & 0 & 0 & 0
\end{pmatrix}
\end{equation*}

A piecewise linear transformation is applied to enhance the contrast of the image using the following function:

\begin{equation*}
T(r) = \begin{cases}
0, & r < 3 \\
2 \cdot (r - 3), & 3 \le r \le 7 \\
15, & r > 7
\end{cases}
\end{equation*}

\begin{enumerate}
\item[(a)] Apply the transformation to each pixel and produce the transformed image matrix.
\item[(b)] Analyze and describe the effect of this transformation on the image contrast by comparing the histogram of the image before and after the transformation.
\end{enumerate}

\vspace{1.0cm}
\section*{Question 2 \quad [CLO1-PL01-C4] \quad 15 Marks}

A grayscale image exhibits the following pixel intensity distribution:

\begin{center}
\begin{tabular}{|c|c|c|c|c|c|c|c|c|c|}
\hline
\textbf{Intensity} & 0 & 1 & 2 & 3 & 4 & 5 & 6 & 7 & \textbf{Total} \\
\hline
\textbf{Frequency} & 50 & 100 & 100 & 150 & 300 & 150 & 100 & 0 & \textbf{1000} \\
\hline
\end{tabular}
\end{center}

\begin{enumerate}
\item[(a)] Perform histogram equalization on the given image.
\item[(b)] Analyze how the visual appearance of the image would change prior to and following histogram equalization.
\item[(c)] Would the resulting histogram be completely uniform? Provide a justification for your answer.
\end{enumerate}

\vspace{1.0cm}
\section*{Question 3 \quad [CLO1-PL01-C4] \quad 15 Marks}

Given the following $5 \times 5$ image patch $\Pi$ and the Laplacian kernel for edge enhancement:

\begin{equation*}
\Pi = \begin{pmatrix}
12 & 12 & 13 & 15 & 14 \\
12 & 14 & 16 & 17 & 15 \\
13 & 15 & 18 & 19 & 16 \\
14 & 16 & 18 & 20 & 17 \\
13 & 14 & 15 & 17 & 16
\end{pmatrix}
\quad
H = \begin{pmatrix}
0 & -1 & 0 \\
-1 & 4 & -1 \\
0 & -1 & 0
\end{pmatrix}
\end{equation*}

\begin{enumerate}
\item[(a)] Apply the Laplacian filter $H$ to the image patch $\Pi$.
\item[(b)] Produce a sharpened image using the Laplacian response.
\item[(c)] Analyze and describe the manner in which high-frequency details are affected by this operation.
\end{enumerate}

\vspace{1.0cm}
\section*{Question 4 \quad [CLO1-PL01-C4] \quad 10 Marks}

Digital image enhancement can be performed in the frequency domain using filters such as Ideal, Butterworth, and Gaussian filters. Analyze and explain the steps involved in frequency domain filtering, addressing the following points:

\begin{enumerate}
\item[(a)] Explain why the Fourier Transform of an image is performed prior to filtering.
\item[(b)] Describe how low-pass and high-pass filters differ in their effect on an image.
\item[(c)] Discuss the advantages and disadvantages of performing filtering in the frequency domain as compared to the spatial domain.
\item[(d)] Provide specific examples of scenarios where frequency domain filtering is preferable to spatial domain filtering.
\end{enumerate}

\vspace{2.0cm}
\noindent\rule{\textwidth}{0.5pt}
\begin{center}
\textit{End of Examination Paper}
\end{center}

\end{document}
